\section{ماژول aa\_svd\_compressor (Activation-Aware SVD Compressor)}
\label{sec:module_aa_svd_compressor}

\subsection{معرفی و انگیزه}
این ماژول در چارچوب پژوهش حاضر پیاده‌سازی و ارزیابی شده و فاقد مقاله مستقل جداگانه در ادبیات تحقیق است.

تقریب رتبه پایین ماتریس‌های وزن ترنسفورمر از طریق تجزیه مقادیر منفرد (\lr{SVD}) در صورتی که صرفاً مبتنی بر ماتریس وزن ایستا باشد، به سرعت با افت شدید دقت در توالی‌های عمیق مواجه می‌شود، زیرا جهت‌گیری مقادیر ویژه ماتریس وزن با جهت‌گیری بیشترین پراکندگی فعال‌سازی‌های پنهان ورودی هم‌راستا نیست. در نتیجه، جهت‌های فرعی وزن که بار اطلاعاتی عظیمی از فعال‌سازی‌های ورودی را حمل می‌کنند هرس شده و مدل دچار رانش شدید بازنمایی (\lr{Representation Drift}) می‌گردد.

فشرده‌ساز \lr{Activation-Aware SVD (AASVD) Compressor} به منظور حل ساختاری این چالش توسعه یافته است. ایده محوری در این ماژول، تخمین برخط ماتریس کوواریانس فعال‌سازی‌های ورودی از طریق میانگین متحرک نمایی (\lr{Streaming EMA Covariance}) و پیش‌شرط‌سازی متقارن ماتریس وزن با جذر کوواریانس است. این فرآیند موجب می‌شود که مؤلفه‌های منفرد وزن‌ها در جهت زیرفضاهای پرتراکم فعال‌سازی همتراز شده و با جایگزینی لایه‌های متراکم با لایه‌های فاکتورگیری‌شده فیزیکی \lr{SVDLinear}، کاهش واقعی در پارامترها و پیچیدگی محاسباتی حاصل شود.

\subsection{مبانی نظری و فرمول‌بندی ریاضی}
در طول فاز کالیبراسیون، ماتریس کوواریانس ورودی $\mathbf{C}_X$ با استفاده از نمونه‌های ورودی و ضریب واپاشی نمایی $\beta \in [0, 1]$ به صورت جاری به‌روزرسانی می‌گردد:
\begin{equation}
\mathbf{C}_X^{(t)} = \beta \mathbf{C}_X^{(t-1)} + (1 - \beta) \frac{1}{B} \sum_{i=1}^B \mathbf{X}_i^\top \mathbf{X}_i
\end{equation}
به منظور تضمین مثبت معین بودن ماتریس، کوواریانس متقارن‌سازی و منظم‌سازی شده و تجزیه مقادیر ویژه بر روی آن اعمال می‌شود:
\begin{equation}
\mathbf{C}_X = \mathbf{Q}_c \mathbf{\Lambda}_c \mathbf{Q}_c^\top, \quad \mathbf{C}_X^{1/2} = \mathbf{Q}_c \mathbf{\Lambda}_c^{1/2} \mathbf{Q}_c^\top, \quad \mathbf{C}_X^{-1/2} = \mathbf{Q}_c \mathbf{\Lambda}_c^{-1/2} \mathbf{Q}_c^\top
\end{equation}
ماتریس وزن اولیه لایه $\mathbf{W} \in \mathbb{R}^{d_{\text{out}} \times d_{\text{in}}}$ با جذر کوواریانس مقیاس‌گذاری متقارن می‌شود:
\begin{equation}
\mathbf{W}_{\text{weighted}} = \mathbf{W} \mathbf{C}_X^{1/2}
\end{equation}
سپس تجزیه مقادیر منفرد بر روی این ماتریس وزن هم‌تراز شده اجرا می‌گردد:
\begin{equation}
\mathbf{W}_{\text{weighted}} = \mathbf{U} \mathbf{\Sigma} \mathbf{V}^\top
\end{equation}
برای رتبه انتخابی $r \le \min(d_{\text{in}}, d_{\text{out}})$، مقادیر منفرد با قاعده جذر متقارن به فاکتورهای چپ و راست تخصیص داده شده و اثر پیش‌شرط‌سازی با ضرب معکوس جذر کوواریانس خنثی می‌گردد:
\begin{equation}
\mathbf{U}_r = \mathbf{U}_{:, :r} \mathbf{\Sigma}_{:r}^{1/2} \in \mathbb{R}^{d_{\text{out}} \times r}, \quad \mathbf{V}_r = \mathbf{\Sigma}_{:r}^{1/2} \mathbf{V}_{:, :r}^\top \mathbf{C}_X^{-1/2} \in \mathbb{R}^{r \times d_{\text{in}}}
\end{equation}
در حین استنتاج، خروجی لایه فشرده با دو پروجکشن متوالی محاسبه می‌شود:
\begin{equation}
\mathbf{y} = (\mathbf{x} \mathbf{V}_r^\top) \mathbf{U}_r^\top + \mathbf{b}
\end{equation}
که در آن تعداد کل پارامترها از $d_{\text{in}} \cdot d_{\text{out}}$ به $r \cdot (d_{\text{in}} + d_{\text{out}})$ تقلیل می‌یابد.

\subsection{معماری و پیاده‌سازی}
کلاس‌های \lr{AASVDCompressor} و \lr{SVDLinear} در فایل \lr{Codes/framework/compression\_engine/aa\_svd\_compressor.py} قرار دارند و با شناسه‌های \lr{aasvd}، \lr{aa\_svd} و \lr{aa\_svd\_compressor} در رجیستری ماژول‌های فشرده‌سازی ثبت شده‌اند. نقاط کلیدی در پیاده‌سازی عبارتند از:
\begin{itemize}
    \item \textbf{جایگزینی ماژولار ساختاری (\lr{SVDLinear Module}):} لایه‌های چگال استاندارد و همچنین لایه‌های \lr{Conv1D} (در مدل \lr{GPT-2}) مستقیماً با ماژول اختصاصی \lr{SVDLinear} جایگزین می‌شوند. این لایه علاوه بر فاکتورهای کم‌مرتبه $\mathbf{U}$ و $\mathbf{V}$، مجهز به آداپتورهای باقیمانده \lr{LoRA} برای فاز بازیابی دقت پس از فشرده‌سازی است.
    \item \textbf{پشتیبانی از هسته همجوش‌یافته سریع (\lr{Fused SVD Kernel}):} کلاس \lr{SVDLinear} امکان فراخوانی هسته‌های بهینه‌سازی‌شده برای ادغام دو ضرب خطی و اعمال بایاس در یک گذر حافظه \lr{SRAM} را داراست.
    \item \textbf{اعتبارسنجی خطای فروبنیوس و مسیر بازگشت ایمن (\lr{Safe Fallback}):} در صورت فراتر رفتن خطای نسبی فروبنیوس از آستانه مجاز (\lr{max\_error=0.50})، عملیات فشرده‌سازی برای لایه مربوطه متوقف شده و وزن‌های متراکم اولیه بازگردانده می‌شوند.
    \item \textbf{پیکربندی هوشمند رتبه فشرده‌سازی (\lr{Dynamic Rank Allocation}):} رتبه لایه‌ها با لحاظ قید کاهش پارامتری $r \cdot (d_{\text{in}} + d_{\text{out}}) < d_{\text{in}} \cdot d_{\text{out}}$ به صورت پویا متناسب با ابعاد هر لایه تنظیم می‌گردد.
\end{itemize}

\subsection{نتایج تجربی}
ارزیابی ماژول \lr{aa\_svd\_compressor} بر روی مدل‌های زبانی و انکودرهای پایه بر اساس داده‌های تجربی تاییدشده در \lr{thesis\_final/master\_verified\_results.json} انجام شده و در جدول \ref{tab:benchmark_aa_svd_compressor} گزارش شده است:

\begin{table}[H]
\centering
\small
\caption{نتایج تجربی و اعتبارسنجی بنچ‌مارک فشرده‌ساز \lr{aa\_svd\_compressor} روی مدل‌های پایه}
\label{tab:benchmark_aa_svd_compressor}
\begin{tabular}{lccccc}
\toprule
\textbf{مدل پایه} \lr{(Base Model)} & \textbf{دقت} \lr{(Accuracy)} & \textbf{شاخص} \lr{(F1-Score)} & \textbf{حافظه} \lr{(Peak VRAM)} & \textbf{توان} \lr{(Throughput)} & \textbf{وضعیت} \\
\midrule
\lr{DistilBERT} & \lr{89.45\%} & \lr{0.8996} & \lr{234.19 MB} & \lr{163,468.26 tok/s} & \lr{PASS} \\
\lr{RoBERTa} & \lr{88.07\%} & \lr{0.8877} & \lr{343.48 MB} & \lr{101,496.89 tok/s} & \lr{PASS} \\
\lr{GPT-2} & \lr{80.73\%} & \lr{0.8243} & \lr{467.74 MB} & \lr{12,376.10 tok/s} & \lr{PASS} \\
\lr{Qwen} & \lr{---} & \lr{---} & \lr{---} & \lr{---} & \lr{PASS} \\
\bottomrule
\end{tabular}
\end{table}

نتایج ثبت‌شده برتری قاطع الگوریتم \lr{AASVD} را در حفظ کیفیت مدل به اثبات می‌رساند. این ماژول موفق به دستیابی به بالاترین دقت‌های ثبت‌شده شامل \lr{89.45\%} در \lr{DistilBERT}، \lr{88.07\%} در \lr{RoBERTa} و \lr{80.73\%} در دیکودر علّی \lr{GPT-2} گردیده است. هم‌زمان، به واسطه سبک‌سازی وزن‌ها، توان پردازشی بالا \lr{163,468.26 tok/s} با مصرف حافظه ناچیز \lr{234.19 MB} به ثبت رسیده که تاییدی بر بهینگی همترازی زیرفضایی است.

\subsection{محدودیت‌ها و تحلیل شکست}
با وجود نتایج قابل قبول، تحلیل مکانیزم‌های داخلی ماژول نشان‌دهنده برخی ویژگی‌های محدودکننده است:
\begin{itemize}
    \item \textbf{سربار فاز تخمین کوواریانس (\lr{Calibration Latency}):} پیش‌شرط‌سازی ماتریس وزن‌ها مستلزم اجرای گذر پیشرو بر روی مجموعه متون کالیبراسیون و محاسبه ضرایب کوواریانس است. در مدل‌های بسیار بزرگ با ابعاد پنهان چند ده‌هزار، محاسبه مقادیر ویژه و جذر ماتریس $\mathbf{C}_X$ نیازمند توان محاسباتی و زمان آماده‌سازی اولیه بالاتری در مقایسه با روش‌های بدون داده است.
    \item \textbf{افت توان نسبی در مدل \lr{GPT-2} (\lr{12,376.10 tok/s}):} تفاوت سرعت پردازش در \lr{GPT-2} نسبت به مدل‌های انکودر به واسطه تقسیم لایه فشرده \lr{c\_attn} به سه بخش مجزا برای سرها و مدیریت تفکیک فاکتورهای کم‌مرتبه در لایه‌های دیکودر متوالی است که نیازمند فراخوانی مکرر هسته‌های خطی متوالی می‌باشد.
    \item \textbf{پایداری عددی مقادیر ویژه کم‌انرژی:} در ابعاد فعال‌سازی بسیار خلوت، مقادیر ویژه انتهایی ماتریس کوواریانس ممکن است به سمت صفر میل کنند؛ در صورت عدم اعمال منظم‌سازی دقیق $\epsilon$ در مرحله معکوس‌سازی جذر کوواریانس، خطر بزرگ‌نمایی خطاهای عددی در فاکتور $\mathbf{V}_r$ وجود خواهد داشت که این ماژول با تنظیم $\epsilon \ge 10^{-8}$ آن را مهار نموده است.
\end{itemize}
